% Le problème du chapitre : l'écriture va au primaire, la lecture suivante à
% une réplique que la modification n'a pas encore atteinte.
\documentclass[border=6pt]{standalone}
\usepackage{sd}
\begin{document}
\begin{tikzpicture}[x=1mm, y=1mm,
    lab/.style={font=\sffamily\small, text=#1, fill=white, inner sep=1.5pt},
    lab/.default=Dink]
  \node[seqhead=Dblue] (u) at (0,0) {\stack{laptop}{Navigateur}};
  \node[seqhead=Dpurple] (p) at (64,0) {\stack{database}{Primaire}};
  \node[seqhead=Dteal] (r) at (128,0) {\stack{database}{Réplique}};
  \seqlife{-92}{u,p,r}
  % retard de réplication, sur la ligne de vie de la réplique
  \fill[DamberBg] (126.6,-26) rectangle (129.4,-62);
  \draw[Damber, line width=0.7pt] (126.6,-26) rectangle (129.4,-62);
  \draw[accol, Damber] (131,-26) -- node[right=6pt, note, align=left, text=Damber, font=\sffamily\footnotesize] {retard de\\réplication :\\la réplique a\\encore l'ancien nom} (131,-62);
  % 1. écriture
  \draw[arr=Dblue] (0,-16) -- node[lab, above=1pt] {modifier : nom = « Léa Martin »} (64,-16);
  \node[num=Dblue] at (0,-16) {1};
  \draw[darr=Dpurple, line width=0.9pt] (64,-26) -- node[lab=Dslate, above=1pt] {OK, enregistré} (0,-26);
  % 2. lecture trop tôt
  \draw[arr=Dblue] (0,-38) -- node[lab, above=1pt, pos=0.25] {afficher le profil} (126.6,-38);
  \node[num=Dblue] at (0,-38) {2};
  \draw[darr=Dred, line width=0.9pt] (126.6,-48) -- node[lab=Dred, above=1pt, pos=0.75] {nom = « Léa » : l'ancienne version} (0,-48);
  % 3. la modification atteint la réplique
  \draw[arr=Dpurple] (64,-62) -- node[lab, above=1pt] {modification répliquée} (126.6,-62);
  \node[num=Dpurple] at (64,-62) {3};
  % 4. lecture après la réplication
  \draw[arr=Dblue] (0,-74) -- node[lab, above=1pt, pos=0.25] {afficher le profil} (128,-74);
  \node[num=Dblue] at (0,-74) {4};
  \draw[darr=Dgreen, line width=0.9pt] (128,-84) -- node[lab=Dgreen, above=1pt, pos=0.75] {nom = « Léa Martin »} (0,-84);
\end{tikzpicture}
\end{document}
